\chapter{Translation and braided exchange}
\label{chap:translation-braided-exchange}

Translated polynomial multiplications commute because their coefficient
algebra is commutative.
For a noncommutative algebra, reversing their order requires an exchange relation.
On the quantum plane, the two monomial degrees determine its scalar coefficient.

Throughout this chapter, $k$ is a field of characteristic zero.
All vector spaces have cohomological degree zero and even parity.
Tensor permutations therefore carry no signs.
All position variables and deformation parameters commute with the
coefficients and have degree zero.
They are formal, with no numerical values or analytic convergence conditions.

\section{Multiplication at two positions}
\label{sec:vertex-polynomial-translation}

For $f,g\in k[t]$, define the position-dependent operator
\[
 Y(f,z)g=f(t+z)g(t)\in k[t,z].
\]
The variable $t$ belongs to the coefficient algebra, while $z$
records the position of the first polynomial.
Applying two operators gives
\[
 \begin{aligned}
 Y(t,z_1)Y(t^2,z_2)1
   &=(t+z_1)(t+z_2)^2\\
   &=t^3+t^2(z_1+2z_2)+t(2z_1z_2+z_2^2)+z_1z_2^2.
 \end{aligned}
\]
Reversing their order gives the same polynomial, because the
coefficient algebra is commutative.
The constant polynomial $1$ acts as the identity operator.

There is also a composition law for positions.
For $f,g,c\in k[t]$, direct substitution gives
\[
 \begin{aligned}
 Y(f,z_0+z_2)Y(g,z_2)c
   &=f(t+z_0+z_2)g(t+z_2)c(t)\\
   &=Y\bigl(f(t+z_0)g(t),z_2\bigr)c\\
   &=Y(Y(f,z_0)g,z_2)c.
 \end{aligned}
\]
In the middle expression, the outer operator substitutes $t+z_2$
for $t$ and leaves the independent variable $z_0$ fixed.
Thus multiplication at two positions equals an iterated multiplication
at their relative displacement.

Polynomial differentiation produces this translation.
For $j\geq0$, the derivative of $t^j$ vanishes after order $j$, and
\[
 \sum_{r\geq0}\frac{z^r}{r!}\partial_t^r(t^j)
   =\sum_{r=0}^j\binom jr t^{j-r}z^r=(t+z)^j.
\]
Linearity proves the same identity for every polynomial.
Differentiation on another coefficient algebra need not have terminating powers.
Its exponential then requires formal series with vector coefficients.

\section{Formal fields and their products}
\label{sec:vertex-formal-carriers}

A pole such as $z^{-1}$ requires negative powers.
Products expanded in different variable orders must then be compared
in a common space of coefficient arrays.
For a $k$-vector space $V$, write
\[
 V[[z]]=\left\{\sum_{j\geq0}v_jz^j:v_j\in V\right\},\qquad
 V((z))=\bigcup_{L\geq0}z^{-L}V[[z]].
\]
A Laurent series has one finite lower bound on its exponents.
Its coefficients need not lie in one finite-dimensional subspace.

The notation $V[[z_1,z_2]]$ denotes arrays with both exponents
nonnegative. Iteration gives $V((z_1))((z_2))$, where the lower
$z_1$ bound may depend on the $z_2$ coefficient.

A formal distribution in $d$ variables is an arbitrary coefficient
array $(v_{n_1,\ldots,n_d})_{(n_1,\ldots,n_d)\in\mathbb Z^d}$.
We write it as $\sum v_{n_1,\ldots,n_d}z_1^{n_1}\cdots z_d^{n_d}$
and compare distributions coefficientwise.
Every iterated Laurent series defines such a distribution.
Multiplication by a polynomial is always defined, since each output
coefficient is a finite sum.
For more general products, coefficientwise finiteness must be proved.

A formal field on $V$ is a linear map $a(z):V\to V((z))$.
Its modes are the linear maps $a_{(n)}:V\to V$ determined by
\[
 a(z)v=\sum_{n\in\mathbb Z}(a_{(n)}v)z^{-n-1}.
\]
The Laurent condition says that, for every $v$, one has
$a_{(n)}v=0$ for all sufficiently large $n$.
This is lower truncation, and its bound may depend on $v$.
A field is regular when it takes values in $V[[z]]$.

The vectors of $V$ are called states.
A state-field map assigns a field $Y(u,z)$ linearly to each $u\in V$.
Thus
\[
 \begin{aligned}
 Y&:V\longrightarrow\operatorname{Hom}_k(V,V((z))),\\
 Y(z)&:V\otimes_kV\longrightarrow V((z)),\qquad
 Y(z)(u\otimes v)=Y(u,z)v
 \end{aligned}
\]
are two forms of the same map.
Here an algebraic tensor is a finite sum of elementary tensors,
subject to bilinearity.
For fixed $u,v,c\in V$, coefficientwise application gives
\[
 Y(u,z_1)Y(v,z_2)c\in V((z_1))((z_2)),\qquad
 Y(v,z_2)Y(u,z_1)c\in V((z_2))((z_1)).
\]
For example, the first product applies $Y(u,z_1)$ to each coefficient
of $Y(v,z_2)c$, whose $z_2$ exponents have a lower bound.

The inverse of $z_1-z_2$ has different expansions in the two
iterated Laurent fields:
\[
 G_{12}=\sum_{n\geq0}z_1^{-n-1}z_2^n
       \in k((z_1))((z_2)),\qquad
 G_{21}=-\sum_{n\geq0}z_2^{-n-1}z_1^n
       \in k((z_2))((z_1)).
\]
Multiplication by $z_1-z_2$ cancels successive terms in each
series and leaves constant term one.
Thus they are the inverses in their respective fields.
As distributions, their difference is
\[
 \delta(z_1,z_2)=G_{12}-G_{21}
      =\sum_{n\in\mathbb Z}z_1^{-n-1}z_2^n.
\]
For $n\leq-1$, putting $j=-n-1$ gives exactly the terms
$z_1^jz_2^{-j-1}$ contributed by $-G_{21}$.

At coefficient $z_1^p z_2^q$, the distribution
$z_1\delta(z_1,z_2)$ has value one when $p+q=0$ and zero otherwise.
The same coefficient rule holds for $z_2\delta(z_1,z_2)$.
Consequently
\[
 (z_1-z_2)\delta(z_1,z_2)=0,
 \qquad \delta(z_1,z_2)\ne0.
\]
The nonzero coefficient at $z_1^{-1}z_2^0$ proves the second assertion.
A nonzero distribution can therefore vanish after multiplication by
a power of $z_1-z_2$.
This permits an exchange condition weaker than equality of the two
ordered products.

The fields $Y(u,z)$ and $Y(v,z)$ are local if some integer
$N\geq0$ satisfies
\begin{equation}\label{eq:vertex-ordinary-locality}
 (z_1-z_2)^N\bigl(Y(u,z_1)Y(v,z_2)c
                  -Y(v,z_2)Y(u,z_1)c\bigr)=0
 \quad\text{for every }c\in V.
\end{equation}
The integer $N$ may depend on $u,v$, but not on $c$.
This equality lies in the space of distributions, so it compares
the coefficients of both iterated products.
For regular fields, both products already belong to $V[[z_1,z_2]]$.

Polynomial translation satisfies $Y(1,z)f=f$, while $Y(f,z)1$ has
constant coefficient $f$.
Write $[z^j]$ for extraction of the coefficient of $z^j$.
A vacuum is a vector $\mathbf1\in V$ with
\[
 Y(\mathbf1,z)v=v,\qquad
 Y(v,z)\mathbf1\in V[[z]],\qquad
 [z^0]Y(v,z)\mathbf1=v.
\]
Differentiating a translated polynomial can act either on its input
or on its position variable.
Use the commutator $[A,B]=AB-BA$, with coefficientwise action on
series, and the formal derivative $\partial_z(z^m)=mz^{m-1}$ for
every integer $m$.
A translation operator is a linear map $D:V\to V$ satisfying
\[
 D\mathbf1=0,\qquad
 [D,Y(v,z)]=Y(Dv,z)=\partial_zY(v,z).
\]

For substitutions involving negative powers, use
\[
 (z_0+z_2)^m
   =\sum_{j\geq0}\binom mj z_0^{m-j}z_2^j,
 \qquad
 \binom mj=\frac{m(m-1)\cdots(m-j+1)}{j!},
\]
where $m\in\mathbb Z$, $j\geq0$, and $\binom m0=1$.
This is expansion in nonnegative powers of $z_2$.
Applied to the first ordered product, it defines
$Y(u,z_0+z_2)Y(v,z_2)c$ in $V((z_0))((z_2))$.
For a fixed $z_2$ coefficient, lower truncation of $Y(v,z_2)c$
bounds the number of input coefficients and binomial terms.
The iterate $Y(Y(u,z_0)v,z_2)c$ belongs to
$V((z_2))((z_0))$ by coefficientwise linearity.

The polynomial composition law equates an ordered product with an iterate.
For Laurent fields, weak associativity requires a polynomial multiple
of their difference to vanish.
For every $u,v,c$, there must be an integer $\ell\geq0$ such that
\begin{equation}\label{eq:vertex-weak-associativity}
 (z_0+z_2)^\ell Y(u,z_0+z_2)Y(v,z_2)c
  =(z_0+z_2)^\ell Y(Y(u,z_0)v,z_2)c.
\end{equation}
The equality compares the resulting distributions with the specified
expansion on the left.
A nonlocal vertex algebra is a vector space with a state-field map
and vacuum satisfying weak associativity.
A vertex algebra here also satisfies ordinary locality.

\section{A derivation produces regular fields}
\label{sec:vertex-differential-algebra}

Let $A$ be a unital associative $k$-algebra with unit $1_A$.
A derivation is a linear map $D:A\to A$ such that
$D(ab)=(Da)b+a(Db)$.
Putting $a=b=1_A$ gives $D1_A=2D1_A$, hence $D1_A=0$.
Define
\begin{equation}\label{eq:vertex-exponential-field}
 E_z(a)=\sum_{j\geq0}\frac{z^j}{j!}D^ja,
 \qquad Y(a,z)b=E_z(a)b,
 \qquad E_z:A\to A[[z]].
\end{equation}
Each coefficient uses one iterate of $D$, so no termination
condition on those iterates is required.

\begin{lemma}\label{lem:vertex-exponential-identities}
For every $a,b\in A$,
\[
 \begin{aligned}
 D^j(ab)&=\sum_{i=0}^j\binom ji(D^ia)(D^{j-i}b),\\
 E_z(ab)&=E_z(a)E_z(b),\qquad E_zE_w(a)=E_{z+w}(a).
 \end{aligned}
\]
The last equality belongs to $A[[z,w]]$, with $E_z$ applied
coefficientwise to $E_w(a)$.
Also $\partial_zE_z(a)=E_z(Da)=D(E_z(a))$.
\end{lemma}

\begin{proof}
The first identity at $j=0$ is multiplication itself.
Apply $D$ to its right side at order $j$.
The derivation rule gives one sum differentiating the first factor
and one differentiating the second.
At an interior index $i$, their combined coefficient is
$\binom j{i-1}+\binom ji=\binom{j+1}i$.
The endpoint coefficients are one, proving induction.

The coefficient of $z^j$ in $E_z(a)E_z(b)$ is
$\sum_{i=0}^j(D^ia)(D^{j-i}b)/(i!(j-i)!)$.
The first identity identifies it with $D^j(ab)/j!$.
The coefficient of $z^pw^q$ in $E_zE_w(a)$ is
$D^{p+q}a/(p!q!)$.
In $E_{z+w}(a)$, only its term of total degree $p+q$ contributes,
and the binomial coefficient gives the same value.
These coefficient calculations also prove that the operations are defined.

Termwise differentiation gives coefficient $D^{j+1}a/j!$ at $z^j$.
This is the coefficient in both $E_z(Da)$ and $D(E_z(a))$.
\end{proof}

\begin{theorem}\label{thm:vertex-differential-algebra}
The formula \eqref{eq:vertex-exponential-field}, with vacuum $1_A$,
makes $A$ a nonlocal vertex algebra with regular fields and
translation operator $D$.
Weak associativity holds with $\ell=0$.
Its fields are ordinarily local for every pair exactly when $A$
is commutative, in which case $N=0$ suffices.
\end{theorem}

\begin{proof}
The definition gives a linear map $Y(a,z):A\to A[[z]]$.
Its modes are
\[
 a_{(-j-1)}b=\frac{D^ja}{j!}b\quad(j\geq0),
 \qquad a_{(n)}b=0\quad(n\geq0),
\]
so lower truncation holds for every input.
Since $D1_A=0$, one has
$Y(1_A,z)b=b$ and $Y(a,z)1_A=E_z(a)$ with constant term $a$.
The coefficient at $z$ in the latter identity is $Da$, so
$a_{(-2)}1_A=Da$.
Thus the field map recovers the translation operator.

For translation, use the derivation rule and the lemma:
\[
 \begin{aligned}
 [D,Y(a,z)]b
   &=D(E_z(a)b)-E_z(a)Db\\
   &=D(E_z(a))b
     =E_z(Da)b=\partial_zY(a,z)b.
 \end{aligned}
\]
The last expression is also $Y(Da,z)b$.
Associativity of $A$ and the two exponential identities give
\begin{equation}\label{eq:vertex-exact-associativity}
 \begin{aligned}
 Y(a,z_0+z_2)Y(b,z_2)c
   &=E_{z_0+z_2}(a)E_{z_2}(b)c\\
   &=E_{z_2}\bigl(E_{z_0}(a)b\bigr)c\\
   &=Y(Y(a,z_0)b,z_2)c.
 \end{aligned}
\end{equation}
All terms belong to $A[[z_0,z_2]]$.
At a fixed total degree, the substitution $z_1=z_0+z_2$
uses finitely many coefficients, proving its validity.

If $A$ is commutative, then $E_{z_1}(a)E_{z_2}(b)c$ equals
$E_{z_2}(b)E_{z_1}(a)c$, giving locality with $N=0$.
Conversely, multiplication by $(z_1-z_2)^N$ is injective on
$A[[z_1,z_2]]$.
Indeed, let $z_2^j$ be the lowest nonzero coefficient of a
nonzero series $F$, viewed in $A[[z_1]][[z_2]]$.
The product's $z_2^j$ coefficient is $z_1^N[z_2^j]F$, which is nonzero.
Cancel the locality factor, apply the equality to $1_A$, and take
the constant coefficient in both variables.
This gives $ab=ba$ for every $a,b$.
\end{proof}

For $A=k[t]$ and $D=\partial_t$, the theorem recovers
$Y(f,z)g=f(t+z)g(t)$ from Section~\ref{sec:vertex-polynomial-translation}.
For the Euler derivation $D=t\partial_t$, it instead gives
\[
 Y(t^m,z)t^r=e^{mz}t^{m+r},\qquad
 e^{mz}=\sum_{j\geq0}\frac{m^jz^j}{j!}.
\]
For instance, $Y(t,z)1=t+tz+tz^2/2+\cdots$ has infinitely many
nonzero position coefficients.

The field target in Theorem~\ref{thm:vertex-differential-algebra}
cannot generally be replaced by $A\otimes_k k((z))$.
For any vector space $U$, the coefficient map
\[
 U\otimes_k k((z))\longrightarrow U((z)),\qquad
 u\otimes f(z)\longmapsto f(z)u
\]
is injective, and its image consists of series whose coefficients
span a finite-dimensional subspace of $U$.
To prove injectivity, express a tensor using independent first factors.
Vanishing of each output coefficient forces every scalar coordinate to vanish.
Conversely, a basis of the finite coefficient span expresses a series
as a finite tensor sum, with Laurent coordinate series.

For a strict inclusion, take the polynomial algebra
$A=k[t_0,t_1,t_2,\ldots]$.
Its monomials use only finitely many variables and form a basis.
Define $Dt_i=t_{i+1}$ and extend by the derivation rule.
Explicitly, on a monomial $\prod_i t_i^{a_i}$ the derivative is
$\sum_{i:a_i>0}a_i t_{i+1}t_i^{a_i-1}\prod_{j\ne i}t_j^{a_j}$.
The sum is finite, and addition of exponents proves the derivation
rule on a product of two monomials.
Then
\[
 Y(t_0,z)1=\sum_{n\geq0}\frac{t_nz^n}{n!}
      \in A[[z]]\setminus\bigl(A\otimes_k k((z))\bigr),
\]
since the basis vectors $t_n$ are linearly independent.

\section{The order of an exchange}
\label{sec:vertex-exchange-order}

For a state-field map on $V$, exchanging two fields preserves
the position attached to each input.
To record both orders, define
\[
 \begin{aligned}
 P_{12}(u\otimes v\otimes c)&=Y(u,z_1)Y(v,z_2)c,\\
 P_{21}(u\otimes v\otimes c)&=Y(u,z_2)Y(v,z_1)c,\\
 \tau(u\otimes v)&=v\otimes u.
 \end{aligned}
\]
Their domains are $V^{\otimes3}$ for the first two maps and
$V^{\otimes2}$ for $\tau$, with tensors over $k$.
The respective product codomains are the iterated spaces in
Section~\ref{sec:vertex-formal-carriers}.
For regular fields, both codomains can be taken as $V[[z_1,z_2]]$.

An unflipped exchange operator is a linear map
\[
 S(x):V\otimes V\longrightarrow V\otimes V\otimes k((x)).
\]
Its name specifies that the identity exchange is $S=1$, the
identity map, rather than the permutation $\tau$.
Each image is a finite tensor sum.
In particular, write
\[
 S(x)(v\otimes u)=\sum_{i=1}^r v_i\otimes u_i\,f_i(x),
 \qquad f_i(x)\in k((x)).
\]
The input here is the reversed ordered pair $v\otimes u$.

For $f\in k((x))$, specify the scalar substitution by
\[
 \iota_{21}f(z_2-z_1)
   =\sum_{j\geq0}\frac{(-z_1)^j}{j!}f^{(j)}(z_2)
   \in k((z_2))[[z_1]].
\]
The derivative is the Laurent derivative defined above.
At a fixed $z_1$ power, the coefficient is one Laurent series
in $z_2$, so this substitution is defined.
On a monomial $f(x)=x^m$, the formula is exactly the binomial
expansion in nonnegative powers of $z_1$.

The corresponding $S$-locality condition is
\begin{equation}\label{eq:vertex-s-locality}
 \begin{aligned}
 &(z_1-z_2)^N Y(u,z_1)Y(v,z_2)c\\
 &\qquad=(z_1-z_2)^N\sum_{i=1}^r
    \iota_{21}f_i(z_2-z_1)Y(v_i,z_2)Y(u_i,z_1)c.
 \end{aligned}
\end{equation}
For every $u,v$, one integer $N\geq0$ must work for all $c$.
The products on the right belong to $V((z_2))((z_1))$.
Their scalar multiplication is defined because the scalar has no
negative $z_1$ powers, while each field product has a lower $z_1$ bound.
A nonlocal vertex algebra satisfying such finite exchange relations
is called a weak quantum vertex algebra.

For regular fields, constant $S$, and $N=0$, the exchange equation is
\begin{equation}\label{eq:vertex-ordered-exchange}
 P_{12}=P_{21}S^{12}\tau^{12}.
\end{equation}
Here a superscript $12$ makes the indicated map act on the first
two factors and leaves the third fixed.
The following diagram displays the two ways to evaluate the product:
\[
 \begin{tikzcd}[column sep=large,row sep=large]
 V^{\otimes3} \arrow[r,"S^{12}\tau^{12}"]
                \arrow[d,"P_{12}"']
   & V^{\otimes3} \arrow[d,"P_{21}"]\\
 V[[z_1,z_2]] \arrow[r,equal]
   & V[[z_1,z_2]].
 \end{tikzcd}
\]
The upper map first sends $u\otimes v$ to $v\otimes u$, then applies $S$.
If its output is $\sum_i v_i\otimes u_i$, the map $P_{21}$
places $v_i$ at $z_2$ and $u_i$ at $z_1$.

The polynomial translation example determines why the reversal is needed.
Take $S=1$, $u=t$, $v=1$, and $c=1$.
Then
\[
 P_{12}(t\otimes1\otimes1)=t+z_1
   =P_{21}(1\otimes t\otimes1),
 \qquad P_{21}(t\otimes1\otimes1)=t+z_2.
\]
Thus $P_{12}=P_{21}$ fails even in this commutative example.
Multiplication of the defect by a locality factor gives
$(z_1-z_2)^{N+1}\ne0$ for every $N\geq0$.

\section{The quantum plane}
\label{sec:vertex-quantum-plane}

Fix $q\in k^\times$, meaning $q\ne0$.
We seek an algebra with generators $a,b$, relation $ba=q\,ab$,
and basis $a^mb^n$ for $m,n\geq0$.
Moving $b^n$ past $a^r$ makes $nr$ exchanges, so the required
multiplication coefficient is $q^{nr}$.

On the vector space
\[
 A_q=\bigoplus_{m,n\geq0}k e_{m,n},
\]
define multiplication bilinearly by
\begin{equation}\label{eq:vertex-plane-multiplication}
 e_{m,n}e_{r,s}=q^{nr}e_{m+r,n+s}.
\end{equation}
The direct sum requires finite support in each element.
For three basis elements with indices $(m,n),(r,s),(p,t)$,
the two parenthesizations have exponents
\[
 nr+(n+s)p=sp+n(r+p).
\]
They have the same basis vector, so multiplication is associative.
The element $e_{0,0}$ is its unit.

Put $a=e_{1,0}$ and $b=e_{0,1}$.
Then $ba=q ab$ and $a^mb^n=e_{m,n}$.
These equations identify $A_q$ with the quotient
$k\langle a,b\rangle/(ba-qab)$, where the free algebra has finite
words as a basis and concatenation as multiplication.
Indeed, replacing $ba$ by $qab$ reduces the number of inverted
pairs in a word until all $a$'s precede all $b$'s.
The resulting ordered words span the quotient, and their images
are the independent basis vectors $e_{m,n}$.
The induced surjective map from the quotient to $A_q$ is therefore injective.

The ordinary derivative in $a$ would not preserve the relation
when $q\ne1$.
An assignment $Da=1$, $Db=0$ would give
$D(ba-qab)=(1-q)b$, a nonzero basis multiple in that case.
The assignment $Da=a$, $Db=b$ sends $ba-qab$ to $2(ba-qab)$.
It therefore preserves the relation and acts on $a^mb^n$ by
multiplication by $m+n$.

For $\alpha=(m,n)\in\mathbb N^2$, write $e_\alpha=e_{m,n}$
and $|\alpha|=m+n$, where $\mathbb N=\{0,1,2,\ldots\}$.
This grading records monomial exponents and contributes no permutation signs.
Define $De_\alpha=|\alpha|e_\alpha$.
Since $|\alpha+\beta|=|\alpha|+|\beta|$, multiplication
\eqref{eq:vertex-plane-multiplication} makes $D$ a derivation.
Theorem~\ref{thm:vertex-differential-algebra} gives
\begin{equation}\label{eq:vertex-plane-field}
 Y(e_{m,n},z)e_{r,s}
   =e^{(m+n)z}q^{nr}e_{m+r,n+s}.
\end{equation}

The generator fields satisfy
\[
 Y(a,z_1)Y(b,z_2)1=e^{z_1+z_2}ab,
 \qquad
 Y(b,z_2)Y(a,z_1)1=q e^{z_1+z_2}ab.
\]
For this reversed generator product, the scalar factor is $q^{-1}$.

Define the scalar function on pairs of degrees
\[
 \rho(\alpha,\beta)=q^{ms-nr},
 \qquad \alpha=(m,n),\quad\beta=(r,s).
\]
The exponent is an integer, so the assumption $q\ne0$ defines
both positive and negative powers.
For degrees $\alpha,\beta,\gamma\in\mathbb N^2$, direct expansion
of the exponent gives
\begin{equation}\label{eq:vertex-bicharacter}
 \begin{aligned}
 \rho(\alpha+\beta,\gamma)
    &=\rho(\alpha,\gamma)\rho(\beta,\gamma),\\
 \rho(\alpha,\beta+\gamma)
    &=\rho(\alpha,\beta)\rho(\alpha,\gamma),\\
 \rho(\beta,\alpha)&=\rho(\alpha,\beta)^{-1}.
 \end{aligned}
\end{equation}
The first two equations define multiplicativity in each argument,
also called the bicharacter property.
Comparison with \eqref{eq:vertex-plane-multiplication} gives
\[
 e_\alpha e_\beta
   =\rho(\beta,\alpha)e_\beta e_\alpha.
\]

Set
\begin{equation}\label{eq:vertex-plane-s}
 S(e_\alpha\otimes e_\beta)
   =\rho(\alpha,\beta)e_\alpha\otimes e_\beta.
\end{equation}
This defines a constant unflipped operator on $A_q\otimes A_q$.
It has an inverse given by replacing $\rho$ by $\rho^{-1}$.
For arbitrary $c\in A_q$, the commutation rule gives
\[
 \begin{aligned}
 Y(e_\alpha,z_1)Y(e_\beta,z_2)c
   &=e^{|\alpha|z_1+|\beta|z_2}e_\alpha e_\beta c\\
   &=\rho(\beta,\alpha)
       Y(e_\beta,z_2)Y(e_\alpha,z_1)c.
 \end{aligned}
\]
Thus $S$-locality holds with $N=0$, including for finite sums
of homogeneous inputs.

For $u=a^2b=e_{2,1}$ and $v=ab^2=e_{1,2}$, the products
are $uv=q e_{3,3}$ and $vu=q^4e_{3,3}$.
Both have Euler degree three, so
\[
 Y(u,z_1)Y(v,z_2)1=q e^{3z_1+3z_2}e_{3,3}
   =q^{-3}Y(v,z_2)Y(u,z_1)1.
\]
Here $S(v\otimes u)=q^{-3}v\otimes u$ supplies exactly that factor.

\section{Compatibility of the exchange operator}
\label{sec:vertex-exchange-axioms}

Two successive exchanges should restore the original ordered pair.
For three inputs, the two ways of reversing their order use
each pair once.
The unflipped operators record those three pair contributions.
Exchanging an input with a product should agree with exchanging it
with each factor before multiplication.

For a map $S:V^{\otimes2}\to V^{\otimes2}$, put
$S^{21}=\tau S\tau$.
On $V^{\otimes3}$, put $S^{12}=S\otimes1$ and $S^{23}=1\otimes S$.
The map $S^{13}$ acts on the first and third factors:
if $S(u\otimes w)=\sum_i u_i\otimes w_i$, then
$S^{13}(u\otimes v\otimes w)=\sum_i u_i\otimes v\otimes w_i$.
All compositions act from right to left.

In the second exchange of a pair, the two inputs have reversed roles.
This accounts for $S^{21}$ in the unitarity equation.
For constant $S$, unitarity and the quantum Yang--Baxter equation require
\begin{equation}\label{eq:vertex-unitarity-yb}
 SS^{21}=S^{21}S=1,\qquad
 S^{12}S^{13}S^{23}=S^{23}S^{13}S^{12}.
\end{equation}
The second equation compares two orders of the three pair operations,
with the tensor positions unchanged by the operator notation.

Exchanging a vacuum factor should leave it fixed.
A constant exchange operator should also commute with translation of an input.
Reversing two positions changes their relative displacement from $z$ to $-z$.
Translation by $z$ returns the resulting state to the original position.

A quantum vertex algebra with constant exchange operator will mean
a weak quantum vertex algebra with such $S$ and the following
compatibilities with its translation operator $D$.
Write $E_z=\sum_{j\geq0}z^jD^j/j!$ and
$Y(-z)(a\otimes b)=Y(a,-z)b$.
The conditions are
\begin{align}
 S(\mathbf1\otimes v)&=\mathbf1\otimes v,
 & [D\otimes1,S]&=0,
 \label{eq:vertex-vacuum-shift}\\
 Y(u,z)v&=E_zY(-z)S(v\otimes u),
 \label{eq:vertex-s-skew}\\
 S(Y(x)\otimes1)&=(Y(x)\otimes1)S^{23}S^{13}.
 \label{eq:vertex-hexagon}
\end{align}
In \eqref{eq:vertex-s-skew}, $E_z$ acts on the output of $Y(-z)$.
Equation~\eqref{eq:vertex-hexagon} is the hexagon identity.
For regular fields, it is an equality from $V^{\otimes3}$ to
$(V\otimes V)[[x]]$.

\begin{theorem}\label{thm:vertex-quantum-plane}
For every $q\in k^\times$, the fields \eqref{eq:vertex-plane-field}
and exchange operator \eqref{eq:vertex-plane-s} give a quantum vertex
algebra with constant exchange operator.
They also satisfy
\[
 S(v\otimes1)=v\otimes1,\qquad
 [1\otimes D,S]=[1\otimes D,S^{-1}]=0,
\]
\[
 S(1\otimes Y(x))=(1\otimes Y(x))S^{12}S^{13}.
\]
All field identities hold in joint power-series spaces.
\end{theorem}

\begin{proof}
Theorem~\ref{thm:vertex-differential-algebra} proves the vacuum,
translation, and weak-associativity conditions.
The calculation after \eqref{eq:vertex-plane-s} proves $S$-locality.
It remains to verify the tensor identities, for which basis tensors
span each algebraic tensor power.

The zero degree has $\rho(0,\alpha)=\rho(\alpha,0)=1$, proving
both vacuum identities for $S$.
On $e_\alpha\otimes e_\beta$, each of $S$, $S^{-1}$, $D\otimes1$,
and $1\otimes D$ acts by a scalar.
They commute, proving both shift equations.
The two unitarity composites have scalar
$\rho(\alpha,\beta)\rho(\beta,\alpha)=1$.
Both sides of the Yang--Baxter equation have scalar
\[
 \rho(\alpha,\beta)\rho(\alpha,\gamma)\rho(\beta,\gamma)
\]
on $e_\alpha\otimes e_\beta\otimes e_\gamma$, proving that equation.

For the hexagon, write $\alpha=(m,n)$ and $\beta=(r,s)$.
The left side of \eqref{eq:vertex-hexagon} sends the basis triple to
\[
 e^{|\alpha|x}q^{nr}\rho(\alpha+\beta,\gamma)
    e_{\alpha+\beta}\otimes e_\gamma.
\]
The right side has the same exponential, multiplication coefficient,
and basis tensor, with exchange coefficient
$\rho(\beta,\gamma)\rho(\alpha,\gamma)$.
These coefficients agree by \eqref{eq:vertex-bicharacter}.
For the second hexagon, the common remaining factor is
$e^{|\beta|x}q^{sp}e_\alpha\otimes e_{\beta+\gamma}$,
where $\gamma=(p,t)$.
Its exchange coefficients are respectively $\rho(\alpha,\beta+\gamma)$
and $\rho(\alpha,\beta)\rho(\alpha,\gamma)$, which also agree.

Finally, take homogeneous $u=e_\alpha$ and $v=e_\beta$.
The exponential identities and the commutation rule give
\[
 \begin{aligned}
 E_zY(-z)S(v\otimes u)
   &=\rho(\beta,\alpha)E_z\bigl(E_{-z}(v)u\bigr)\\
   &=\rho(\beta,\alpha)vE_z(u)
     =E_z(u)v=Y(u,z)v.
 \end{aligned}
\]
The identity $E_zE_{-z}=1$ follows from
Lemma~\ref{lem:vertex-exponential-identities} by the total-degree
substitution $w=-z$.
Every series in this calculation is regular, so each coefficient
uses finitely many terms.
Bilinearity proves skew symmetry for arbitrary inputs.
\end{proof}

At $q=1$ the multiplication is commutative and $S=1$.
Roots of unity require no change in any proof, since only
multiplicativity and invertibility of powers of $q$ were used.

\section{Unflipped and flipped equations}
\label{sec:vertex-flipped-operators}

Composing $S$ with $\tau$ in the two orders gives the maps
\[
 \begin{aligned}
 B=\tau S,\qquad
 B(e_\alpha\otimes e_\beta)
   &=\rho(\alpha,\beta)e_\beta\otimes e_\alpha,\\
 C=S\tau,\qquad
 C(e_\alpha\otimes e_\beta)
   &=\rho(\beta,\alpha)e_\beta\otimes e_\alpha.
 \end{aligned}
\]
The ordered field equation is $P_{12}=P_{21}C^{12}$.
Both flipped maps become $\tau$ when $q=1$, whereas $S$ becomes
the identity.

The maps $B$ and $C$ coincide exactly when $q^2=1$.
Indeed, their coefficients on $a\otimes b$ are $q$ and $q^{-1}$,
so equality forces $q^2=1$.
Conversely, if $q^2=1$, then $q^j=q^{-j}$ for every integer $j$,
and their displayed basis formulas agree.

Applying $B$ twice gives scalar
$\rho(\alpha,\beta)\rho(\beta,\alpha)=1$ and restores the pair.
Hence $B^2=1$, and the same calculation gives $C^2=1$.
The flipped braid equation is
\[
 B^{12}B^{23}B^{12}=B^{23}B^{12}B^{23}.
\]
On a triple of degrees $(\alpha,\beta,\gamma)$, the left side
visits $(\beta,\alpha,\gamma)$, $(\beta,\gamma,\alpha)$,
and $(\gamma,\beta,\alpha)$.
Its successive factors are $\rho(\alpha,\beta)$,
$\rho(\alpha,\gamma)$, and $\rho(\beta,\gamma)$.

The right side visits $(\alpha,\gamma,\beta)$,
$(\gamma,\alpha,\beta)$, and $(\gamma,\beta,\alpha)$.
Its factors are the same three scalars in reverse order.
This proves the braid equation on all basis triples.
Replacing $\rho$ by $\rho^{-1}$ proves the same equation for $C$.

These flipped identities differ from unflipped unitarity.
Define $B^{21}=\tau B\tau$ as for $S$.
For the generators, $B(a\otimes b)=q b\otimes a$ and
$B^{21}(b\otimes a)=q a\otimes b$.
Therefore
\[
 B^{21}B(a\otimes b)=q^2a\otimes b,
\]
which differs from $a\otimes b$ when $q^2\ne1$.
Thus $B$ satisfies $B^2=1$ but fails the unflipped unitarity
equation at those parameters.

The unflipped Yang--Baxter equation also distinguishes these operators.
On $a\otimes b\otimes b$, its two sides with $B$ are
\[
 \begin{aligned}
 B^{12}B^{13}B^{23}(a\otimes b\otimes b)
   &=q\,b\otimes b\otimes a,\\
 B^{23}B^{13}B^{12}(a\otimes b\otimes b)
   &=q^2b\otimes b\otimes a.
 \end{aligned}
\]
For the first composite, $B^{23}$ fixes the two $b$'s,
$B^{13}$ exchanges $a,b$ with factor $q$, and $B^{12}$ fixes the two $b$'s.
For the second, $B^{12}$ and $B^{23}$ each exchange $a,b$
with factor $q$, while the intervening $B^{13}$ fixes the two $b$'s.
The outputs differ for every $q\ne1$, including $q=-1$.

Unitarity also converts the placement of the constant operator in
the field equation.
Since $S\tau S=\tau$, composing
\eqref{eq:vertex-ordered-exchange} with $S^{12}$ gives
\begin{equation}\label{eq:vertex-other-placement}
 P_{12}S^{12}=P_{21}\tau^{12}.
\end{equation}
Conversely, $\tau S^{-1}=S\tau$ converts this equation back.

\section{Exchange depending on displacement}
\label{sec:vertex-spectral-dictionary}

When exchange depends on relative displacement, the three pair
separations are $x$, $x+y$, and $y$.
For example, the positions $x+y,y,0$ have exactly those differences.
This determines the arguments in the spectral Yang--Baxter equation.
Let $V$ have a state-field map, vacuum, and translation operator $D$, and
let $S(x):V\otimes V\to V\otimes V\otimes k((x))$.
The spectral variable $x$ records the relative displacement.

Extend $S(x)$ linearly over $k((x))$ when forming a composite
in that variable, and set $S^{21}(x)=\tau S(x)\tau$.
The corresponding unflipped conditions are
\begin{align}
 S^{21}(-x)S(x)&=S(x)S^{21}(-x)=1,
 \label{eq:vertex-spectral-unitarity}\\
 S^{12}(x)S^{13}(x+y)S^{23}(y)
   &=S^{23}(y)S^{13}(x+y)S^{12}(x).
 \label{eq:vertex-spectral-yb}
\end{align}
The first equation acts on $V\otimes V\otimes k((x))$.
In the second, expand every scalar factor $f(x+y)$ in
nonnegative powers of $y$:
\[
 f(x+y)=\sum_{j\geq0}\frac{y^j}{j!}f^{(j)}(x)
       \in k((x))[[y]].
\]
Both composites then map $V^{\otimes3}$ to
$V^{\otimes3}\otimes k((x))((y))$.
Each application of $S$ has finitely many tensor terms, and
the scalar products belong to this common iterated Laurent field.

The vacuum and translation conditions become
\[
 \begin{aligned}
 S(x)(\mathbf1\otimes v)&=\mathbf1\otimes v,
 &S(x)(v\otimes\mathbf1)&=v\otimes\mathbf1,\\
 [D\otimes1,S(x)]&=-\partial_xS(x).
 \end{aligned}
\]
Here differentiation acts on each scalar Laurent coefficient of $S$.
For an independent formal variable $t$, the shift condition is
equivalent to the translation rule
\[
 (E_t\otimes1)S(x)(E_{-t}\otimes1)=S(x-t),
 \qquad E_t=\sum_{j\geq0}\frac{t^jD^j}{j!}.
\]
Put $H=D\otimes1$ and define $\operatorname{ad}_H(T)=[H,T]$.
Expansion of the exponentials gives coefficient
\[
 \sum_{i=0}^j\frac{(-1)^{j-i}H^iS(x)H^{j-i}}{i!(j-i)!}
   =\frac{\operatorname{ad}_H^j(S(x))}{j!}.
\]
The equality follows by induction from $\operatorname{ad}_H(T)=HT-TH$
and the binomial coefficient identity.

The shift equation identifies this coefficient with
$(-1)^j\partial_x^jS(x)/j!$, by induction since scalar differentiation
commutes with $H$.
This is the Taylor coefficient of $S(x-t)$.
The exponential products have finite coefficient sums, so this is
an identity in $(V\otimes V\otimes k((x)))[[t]]$ on each input.
Conversely, the coefficient of $t$ in the translation rule gives
the shift condition.

Reversing the displacement and translating the output gives the
spectral skew-symmetry condition
\[
 Y(u,z)v=E_zY(-z)S(-z)(v\otimes u),\qquad
 E_z=\sum_{j\geq0}\frac{z^jD^j}{j!}.
\]
The right side belongs to $V((z))$ because the finite tensor
image of $S$ gives a common Laurent lower bound before applying $E_z$.
Each output coefficient of $E_z$ then uses only finitely many derivatives.

Compatibility with multiplication in either factor has the forms
\begin{align}
 S(w)(Y(x)\otimes1)
   &=(Y(x)\otimes1)S^{23}(w)S^{13}(x+w),
 \label{eq:vertex-spectral-hexagon-one}\\
 S(w)(1\otimes Y(x))
   &=(1\otimes Y(x))S^{12}(w-x)S^{13}(w).
 \label{eq:vertex-spectral-hexagon-two}
\end{align}
Expand $f(w\pm x)$ as
$\sum_{j\geq0}(\pm x)^jf^{(j)}(w)/j!$.
Thus both equations compare maps from $V^{\otimes3}$ to
$(V\otimes V)((w))((x))$.
On the left, applying $Y(x)$ first gives a Laurent lower
bound in $x$, and $S(w)$ acts on each coefficient.
On the right, the finite tensor images of $S$ give finitely
many field inputs with a common Laurent lower bound in $x$.
The shifted scalar series have nonnegative $x$ powers, so each
coefficient uses a finite convolution in that variable.

These equations impose conditions on the displacement-dependent operator.
The constant operator of Theorem~\ref{thm:vertex-quantum-plane}
satisfies them: its derivative vanishes, and each shifted value is
the same operator already used in the tensor proofs.
For this operator, the hexagon outputs lie in the smaller space
$(A_q\otimes A_q)[[x]]$.

Unitarity also gives an algebraic placement identity for a formal
displacement variable $t$:
\begin{equation}\label{eq:vertex-spectral-placement}
 \tau S(t)^{-1}=S(-t)\tau.
\end{equation}
Indeed, \eqref{eq:vertex-spectral-unitarity} gives
$S(t)^{-1}=\tau S(-t)\tau$, and left multiplication by $\tau$
proves \eqref{eq:vertex-spectral-placement}.
For the flipped operators $B(t)=\tau S(t)$ and $C(t)=S(t)\tau$,
this gives $B(-t)B(t)=C(-t)C(t)=1$.
For example, $B(-t)B(t)=S^{21}(-t)S(t)$, while the identity
for $C$ follows from $S(-t)S^{21}(t)=1$.

In a field product, the two displacement choices are $t=z_1-z_2$
and $-t=z_2-z_1$.
The first uses the scalar expansion in nonnegative powers of $z_2$,
and the second uses $\iota_{21}$ from
Section~\ref{sec:vertex-exchange-order}.
Equation~\eqref{eq:vertex-spectral-placement} concerns Laurent-valued
tensor maps before these expansions.
Using it inside a locality equation still requires the indicated
expansion spaces and locality factors.
For constant $S$, both scalar expansions agree and give exactly
the conversion \eqref{eq:vertex-other-placement}.

\section{The delta identity and regular Jacobi identity}
\label{sec:vertex-delta}

For regular fields, exact associativity and exact exchange can be
combined into one distribution identity.
The required scalar identity has three terms, corresponding to the
two ordered products and their iterate.
Write
$\delta(t)=\sum_{n\in\mathbb Z}t^n$ and define
\[
 \begin{aligned}
 \Delta_1&=z_0^{-1}\delta\left(\frac{z_1-z_2}{z_0}\right),\\
 \Delta_2&=z_0^{-1}\delta\left(\frac{z_2-z_1}{-z_0}\right),\\
 \Delta_3&=z_2^{-1}\delta\left(\frac{z_1-z_0}{z_2}\right).
 \end{aligned}
\]
These symbols mean binomial expansions of their numerators in
nonnegative powers of $z_2,z_1,z_0$, respectively.
They define scalar distributions in $z_0,z_1,z_2$.
Every monomial appearing in any $\Delta_i$ has total degree $-1$.

\begin{lemma}\label{lem:vertex-delta-identity}
With these expansion conventions,
\begin{equation}\label{eq:vertex-delta-identity}
 \Delta_1-\Delta_2=\Delta_3.
\end{equation}
For every $F\in V[[z_1,z_2]]$, all products in
\[
 \Delta_3F(z_1,z_2)=\Delta_3F(z_0+z_2,z_2)
\]
are coefficientwise defined, and the equality holds.
Multiplication of \eqref{eq:vertex-delta-identity} by $F(z_1,z_2)$
is also coefficientwise defined.
\end{lemma}

\begin{proof}
Consider the coefficient of $z_0^{p_0} z_1^{p_1} z_2^{p_2}$, where $p_0,p_1,p_2$
are integers. All three coefficients vanish unless $p_0+p_1+p_2=-1$.
Under that relation, direct binomial expansion gives the coefficients
\[
 \begin{array}{c|c}
 \Delta_1&
   \begin{cases}(-1)^{p_2}\binom{p_1+p_2}{p_2}&p_2\geq0,\\0&p_2<0,
   \end{cases}\\[4pt]
 \Delta_2&
   \begin{cases}(-1)^{p_2}\binom{p_1+p_2}{p_1}&p_1\geq0,\\0&p_1<0,
   \end{cases}\\[4pt]
 \Delta_3&
   \begin{cases}(-1)^{p_0}\binom{-p_2-1}{p_0}&p_0\geq0,\\0&p_0<0.
   \end{cases}
 \end{array}
\]
For the second row, its sign before simplification is
$(-1)^{-p_0-1+p_1}$, equal to $(-1)^{p_2}$ by the total-degree relation.
Each entry arises from one term of the corresponding binomial expansion.

If $p_0<0$, then $p_1+p_2\geq0$.
When $p_1,p_2\geq0$, the first two binomial coefficients agree.
If $p_2<0$, then $p_1>p_1+p_2$, so the second binomial coefficient is zero
and the first row is absent.
The case $p_1<0$ has the same conclusion with the roles interchanged.
Thus the first two rows cancel and the third is zero.

Suppose $p_0\geq0$, so $p_1+p_2=-p_0-1<0$.
If $p_1,p_2<0$, the first two rows vanish and
$p_0>-p_2-1\geq0$ makes the third vanish.
If $p_2\geq0$, then $p_1<0$, and the first row equals
$\binom{p_0+p_2}{p_2}$, equal to the third row $\binom{p_0+p_2}{p_0}$.
Here the identity $\binom{-a-1}{j}=(-1)^j\binom{a+j}{j}$
for $a,j\geq0$ follows by extracting a minus sign from each numerator factor.

If $p_1\geq0$, then $p_2<0$ and only the second row occurs on the left.
Its negative equals
$(-1)^{p_2+p_1+1}\binom{p_0+p_1}{p_1}=(-1)^{p_0}\binom{p_0+p_1}{p_1}$.
The third row is $(-1)^{p_0}\binom{p_0+p_1}{p_0}$, which is the same number.
These cases exhaust all triples and prove \eqref{eq:vertex-delta-identity}.

Multiplication of $\Delta_3$ by $z_1-z_0-z_2$ gives zero.
Indeed, Pascal's coefficient identity gives
$(z_1-z_0)(z_1-z_0)^n=(z_1-z_0)^{n+1}$ for every integer $n$.
In the defining sum for $\Delta_3$, reindexing $n+1$ then identifies
multiplication by $z_1-z_0$ with multiplication by $z_2$.
All these operations involve polynomial multipliers, so every coefficient is finite.
It follows by repeated polynomial multiplication that
$\Delta_3z_1^az_2^b=\Delta_3(z_0+z_2)^az_2^b$ for $a,b\geq0$.

For a fixed output monomial of total degree $d$, multiplication by
any $\Delta_i$ uses only the total-degree $d+1$ part of $F$.
That part is a finite sum, or zero if $d+1<0$.
The substituted series $F(z_0+z_2,z_2)$ has the same property,
because substitution preserves total degree.
Thus all claimed products are defined, and the polynomial identity
extends to every regular $F$ degree by degree.
\end{proof}

For constant $S$, write $S(v\otimes u)=\sum_i v_i\otimes u_i$,
absorbing any scalar coefficients into the vectors.
The $S$-Jacobi identity in these conventions is the distribution equality
\begin{equation}\label{eq:vertex-s-jacobi}
 \begin{aligned}
 &\Delta_1Y(u,z_1)Y(v,z_2)c
   -\Delta_2\sum_iY(v_i,z_2)Y(u_i,z_1)c\\
 &\qquad=\Delta_3Y(Y(u,z_0)v,z_2)c.
 \end{aligned}
\end{equation}
The quantum-plane fields satisfy this identity directly.
Indeed, put $F(z_1,z_2)=Y(u,z_1)Y(v,z_2)c$.
Exact $S$-locality identifies the second product with this same regular $F$.
Lemma~\ref{lem:vertex-delta-identity} then makes the left side
$\Delta_3F(z_0+z_2,z_2)$.
Exact associativity \eqref{eq:vertex-exact-associativity} identifies
the substituted series with the iterate on the right.

For a commutative differential algebra, the same proof with $S=1$
gives the ordinary Jacobi identity.
Regularity makes each coefficient sum finite.
This argument does not establish the corresponding finiteness for singular fields.

\section{A formal deformation of the quantum plane}
\label{sec:vertex-hadic}

The relation $ba=q\,ab$ becomes commutative at $q=1$.
A formal deformation about this value uses $q=e^h$, where
$h$ is a formal parameter and
\[
 e^h=1+h+\frac{h^2}{2}+\frac{h^3}{6}+\cdots.
\]
Its desired generator relation $ba=e^h ab$ has infinitely many
parameter coefficients.
We also impose identity reduction modulo $h$ on the deformed
exchange operator.

Set $R=k[[h]]$ and let $W=\bigoplus_{m,n\geq0}k e_{m,n}$.
Define
\[
 V=W[[h]]=\left\{\sum_{j\geq0}h^jw_j:w_j\in W\right\}.
\]
A module of this form is called topologically free over $R$.
The submodules $h^NV$, for $N\geq1$, form basic neighborhoods of zero.
A sequence converges when each coefficient of $h$ eventually equals
the corresponding coefficient of its limit.

Every $h$ coefficient in $V$ has finite monomial support.
That support may grow with the coefficient order.
Put $R_N=R/h^NR$ and $V_N=V/h^NV$ for $N\geq1$.
Truncation identifies $V_N$ with $W\otimes_kR_N$.
Define the completed tensor product over $R$ by
\begin{equation}\label{eq:vertex-hadic-tensor}
 V\widehat\otimes_RV
   =\varprojlim_N(V_N\otimes_{R_N}V_N)
   \cong(W\otimes_kW)[[h]].
\end{equation}
The transition maps reduce coefficients from order $N+1$ to order $N$.
Each inverse limit here consists of families whose coefficients agree
under all these reductions.

At level $N$, the isomorphism sends
$(u\otimes r)\otimes(v\otimes s)$ to $(u\otimes v)\otimes rs$.
Its inverse sends $(u\otimes v)\otimes r$ to
$(u\otimes1)\otimes(v\otimes r)$.
Scalar balancing proves both formulas well-defined and mutually inverse.
Compatible truncated coefficients then prove
\eqref{eq:vertex-hadic-tensor}.
This completes with respect to powers of $h$ over the base $R$.
It specifies a different filtration and scalar base from the
two-variable tensor completion in Section~\ref{sec:cc-tensor}.

For $r\geq2$, the same construction defines
\[
 V^{\widehat\otimes_R r}
   =\varprojlim_N V_N^{\otimes_{R_N}r}
   \cong W^{\otimes_k r}[[h]].
\]
At level $N$, the isomorphism multiplies the $r$ scalar factors
and retains the $r$ vector factors.
Its inverse puts the scalar on the last vector factor and
puts $1$ on every preceding factor.
Balancing and compatible coefficients prove the isomorphism as for $r=2$.
Give these completed tensor powers the topology with basic neighborhoods
$h^NV^{\widehat\otimes_R r}$.

Define the multiplication and exchange on constant basis elements by
\begin{equation}\label{eq:vertex-hadic-operations}
 \begin{aligned}
 e_{m,n}*_h e_{r,s}&=e^{hnr}e_{m+r,n+s},\\
 De_{m,n}&=(m+n)e_{m,n},\\
 S_h(e_{m,n}\otimes e_{r,s})
   &=e^{h(ms-nr)}e_{m,n}\otimes e_{r,s}.
 \end{aligned}
\end{equation}
Here $e^{hc}=\sum_{j\geq0}c^jh^j/j!$ for every $c\in k$.

In particular, $a*_h b=ab=e_{1,1}$ and $b*_h a=e^h ab$.
Their difference is
\[
 b*_h a-a*_h b
   =h ab+\frac{h^2}{2}ab+\frac{h^3}{6}ab+\cdots.
\]
Also $S_h(a\otimes b)=e^h a\otimes b$, with the same
first-order scalar coefficient.

Extend the maps $R$-linearly and continuously.
For multiplication, the coefficient of $h^d$ on general inputs
$\sum_i h^iu_i$ and $\sum_j h^jv_j$ is
\[
 \sum_{i+j+\ell=d}\mu_\ell(u_i,v_j),\qquad
 \mu_\ell(e_{m,n},e_{r,s})=\frac{(nr)^\ell}{\ell!}e_{m+r,n+s},
\]
where $i,j,\ell\geq0$ and $\mu_\ell$ is bilinear.
There are finitely many such triples, and each $u_i,v_j$ has
finite support, so the coefficient belongs to $W$.
The same finite-convolution argument defines $S_h$ on
$(W\otimes W)[[h]]$ and defines $D$ on $V$.

All these maps preserve parameter order, proving continuity.
Their formulas on polynomial truncations determine them uniquely, because
$\bigcap_Nh^NV=0$ and likewise for the completed tensor product.
For $c,d\in k$, the identity $e^{hc}e^{hd}=e^{h(c+d)}$ follows at coefficient $h^j$
from the finite binomial sum for $(c+d)^j/j!$.
Consequently the associativity and bicharacter proofs for $q$ apply
coefficientwise with $q=e^h$.
The inverse of $S_h$ replaces $h(ms-nr)$ by its negative.
In particular, $S_h=1+O(h)$ means $(S_h-1)(V\widehat\otimes_RV)
\subset h(V\widehat\otimes_RV)$.

Define the field-output space for this topology by
\[
 V((z))_h=\varprojlim_N V_N((z)).
\]
Equivalently, it consists of distributions $\sum_{j\in\mathbb Z}v_jz^j$
with $v_j\in V$ and $v_j\to0$ in the $h$-adic topology as
$j\to-\infty$.
To prove the description, a compatible family determines each $v_j$
by its compatible truncated coefficients.
At every level $N$, its finite Laurent lower bound makes
$v_j\in h^NV$ for sufficiently negative $j$.
Conversely, these conditions give a Laurent series modulo every $h^N$
and hence a compatible family.

Give $V[[z]]$ the $h$-adic topology, with basic neighborhoods
$h^NV[[z]]$.
Coefficient reduction identifies it with
\[
 V[[z]]=\varprojlim_N V_N[[z]].
\]
Indeed, compatible reductions determine each position coefficient in $V$,
and an arbitrary family of those coefficients gives all the reductions.
Use the same topology on $V[[z_1,\ldots,z_r]]$ for each finite $r$.
Give $V((z))_h$ the topology whose basic neighborhoods are kernels
of its projections to $V_N((z))$.
The inclusion $V[[z]]\to V((z))_h$ preserves the coefficient
reductions and is therefore continuous.

The field formula
\[
 Y_h(u,z)v=(e^{zD}u)*_h v
\]
defines a continuous $R$-linear map
$Y_h(z):V\widehat\otimes_RV\to V[[z]]\subset V((z))_h$.
For a fixed $h^dz^r$ coefficient, the multiplication convolution
has finitely many terms and the exponential contributes $D^r/r!$.
The map $Y_h(z)$ sends $h^NV^{\widehat\otimes_R2}$ into
$h^NV[[z]]$, proving continuity with a bound uniform in the position coefficient.
The scalar target for the exchange map can be taken as
$(W\otimes W\otimes k((x)))[[h]]$.
Give it the topology of its powers of $h$ as well.
The constructed $S_h$ lands in its smaller constant-in-$x$ subspace
$(W\otimes W)[[h]]$.

Define the two maps from $V^{\widehat\otimes_R3}$ to
$V[[z_1,z_2]]$ by
\[
 \begin{aligned}
 P_{12,h}(u\otimes v\otimes c)&=Y_h(u,z_1)Y_h(v,z_2)c,\\
 P_{21,h}(u\otimes v\otimes c)&=Y_h(u,z_2)Y_h(v,z_1)c.
 \end{aligned}
\]
Their uniform parameter bounds extend these formulas continuously
to the completed tensor cube.
For unitary $S_h$, exact $S$-locality has the equivalent placement
$P_{12,h}S_h^{12}=P_{21,h}\tau^{12}$ by
\eqref{eq:vertex-other-placement}.
The locality requirement for the formal deformation allows a polynomial
factor after each parameter reduction:
\begin{equation}\label{eq:vertex-hadic-locality}
 (z_1-z_2)^N P_{12,h}S_h^{12}
   \equiv (z_1-z_2)^N P_{21,h}\tau^{12}\pmod{h^M}.
\end{equation}
For each $u,v$ and each $M\geq1$, some $N\geq0$ must make
the equality hold on $u\otimes v\otimes c$ for all $c$.
The remaining requirements are the vacuum and translation equations,
unitarity, the quantum Yang--Baxter equation, shift, and the first hexagon,
interpreted after every parameter reduction.
Together with $S_h=1+O(h)$, these define the formal deformation
convention used here.

\begin{theorem}\label{thm:vertex-hadic-plane}
The operations \eqref{eq:vertex-hadic-operations} and fields $Y_h$
satisfy these formal deformation requirements.
They have exact weak associativity, and
\eqref{eq:vertex-hadic-locality} holds with $N=0$ for every $M$.
Reduction modulo $h$ gives the commutative differential algebra
\[
 V/hV=k[a,b],\qquad D=a\partial_a+b\partial_b,
 \qquad S_h\bmod h=1.
\]
\end{theorem}

\begin{proof}
The scalar exponential law preserves the multiplication and bicharacter
identities under $q^j\mapsto e^{hj}$.
The vacuum is $e_{0,0}$, and $D$ is a derivation,
since both depend only on additivity of the degrees.
The coefficient proof of
Theorem~\ref{thm:vertex-differential-algebra} therefore proves both
translation equations and exact weak associativity over $R$.

For $\alpha=(m,n)$ and $\beta=(r,s)$, put
$\rho_h(\alpha,\beta)=e^{h(ms-nr)}$.
For $c\in V$, the locality calculation is
\[
 \begin{aligned}
 P_{12,h}S_h^{12}(e_\alpha\otimes e_\beta\otimes c)
  &=e^{|\alpha|z_1+|\beta|z_2}
       \rho_h(\alpha,\beta)(e_\alpha*_h e_\beta)*_h c\\
  &=e^{|\alpha|z_1+|\beta|z_2}
       (e_\beta*_h e_\alpha)*_h c\\
  &=P_{21,h}\tau^{12}(e_\alpha\otimes e_\beta\otimes c),
 \end{aligned}
\]
This is an equality in $V[[z_1,z_2]]$.
The vacuum, shift, unitarity, Yang--Baxter, skew-symmetry, and both
hexagon calculations use exactly the same scalar identities.
At any fixed parameter order, arbitrary completed inputs have only
finitely many relevant basis coefficients.
Multilinearity and continuity therefore extend all these equalities
to completed inputs and prove every required reduction.

Modulo $h$, the multiplication coefficient $e^{hnr}$ is one and
the exchange coefficient is one.
Thus the ordered basis monomials multiply as those of $k[a,b]$.
The degree formula for $D$ identifies it with
$a\partial_a+b\partial_b$, giving the asserted reduction.
\end{proof}

The regular Jacobi calculation of Section~\ref{sec:vertex-delta}
also holds coefficientwise in $h$.
At each parameter coefficient, its total-degree argument again uses
finitely many position coefficients.

\begin{proposition}\label{prop:vertex-fixed-q-obstruction}
Fix $q_0\in k^\times$ with $q_0\ne1$.
Extend the multiplication and Euler fields of $A_{q_0}$ to
$A_{q_0}[[h]]$ without changing their coefficients.
No exchange map $S_h=1+O(h)$ can make these fields satisfy
\eqref{eq:vertex-hadic-locality}.
\end{proposition}

\begin{proof}
At $M=1$, reduction of that locality equation gives ordinary
locality for the fields of $A_{q_0}$.
For $u=a$, $v=b$, and $c=1$, its commutator is
\[
 [Y(a,z_1),Y(b,z_2)]1
   =(1-q_0)e^{z_1+z_2}ab\ne0.
\]
The coefficient at $z_1^0z_2^0$ is the nonzero vector
$(1-q_0)e_{1,1}$.
Injectivity of multiplication by $(z_1-z_2)^N$ on regular series,
proved in Theorem~\ref{thm:vertex-differential-algebra}, excludes every $N$.
This contradicts the required reduction independently of the proposed
higher coefficients of $S_h$.
\end{proof}

\section{Nondegeneracy of field evaluation}
\label{sec:vertex-nondegeneracy}

An injective evaluation map would allow equality of field products
on the vacuum to imply equality of their tensor inputs.
Let $K_n=k((z_1))\cdots((z_n))$ be the scalar iterated Laurent field.
For a nonlocal vertex algebra $V$, define
\[
 \begin{aligned}
 Z_n:V^{\otimes n}\otimes_k K_n
   &\longrightarrow V((z_1))\cdots((z_n)),\\
 (u_1\otimes\cdots\otimes u_n)\otimes f
   &\longmapsto fY(u_1,z_1)\cdots Y(u_n,z_n)\mathbf1.
 \end{aligned}
\]
Successive coefficientwise field action gives the stated iterated target.
Scalar multiplication is defined there by iterated Laurent convolution,
whose lower bounds make each coefficient sum finite at each iteration.
Nondegeneracy means that $Z_n$ is injective for every $n\geq1$.

For $V=A_q$ with the Euler fields, put
\[
 \xi=(a\otimes b-q^{-1}b\otimes a)\otimes1
       \in A_q^{\otimes2}\otimes_kK_2.
\]
Let $\lambda_a,\lambda_b:A_q\to k$ extract the coefficients of
$a$ and $b$, respectively.
The scalar extension of $\lambda_a\otimes\lambda_b$ sends $\xi$ to
$1\in K_2$, proving $\xi\ne0$.
The field calculation is
\[
 \begin{aligned}
 Z_2((a\otimes b)\otimes1)&=e^{z_1+z_2}ab,\\
 Z_2((b\otimes a)\otimes1)&=e^{z_1+z_2}ba
                         =q e^{z_1+z_2}ab.
 \end{aligned}
\]
Hence $Z_2(\xi)=0$, proving failure of nondegeneracy for every
$q\in k^\times$, including $q=1$.

For the completed example, define
\[
 \widehat Z_2:V\widehat\otimes_RV\longrightarrow V[[z_1,z_2]],
 \qquad u\otimes v\longmapsto Y_h(u,z_1)Y_h(v,z_2)1.
\]
The coefficient bounds of Section~\ref{sec:vertex-hadic} extend this
map continuously from elementary tensors.
It has the relation $a\otimes b-e^{-h}b\otimes a$ in its domain.
This tensor is nonzero because its reduction modulo $h$ is the difference
of two distinct tensor basis vectors.
Its image under $\widehat Z_2$ vanishes by $b*_h a=e^h ab$.

These evaluation maps have nonzero kernels, so equality of their
outputs does not determine equality of their tensor inputs.
The tensor identities for $S$ and $S_h$ follow from the explicit
degree calculations.
